% Chapter 1 -- the functional model and its solution.
\chapter{The linear model and its solution}\label{ch:model}

\section{Gauss--Markov model}

The observations $\vc l\in\mathbb R^n$ are modelled as
\begin{equation}\label{eq:gm}
\vc l = \mat A\vc x + \vc e,\qquad \E\{\vc e\}=\vc 0,\qquad \mathrm{C}\{\vc e\}=\sigma_0^2\mat\Sigma ,
\end{equation}
with the design matrix $\mat A$ ($n\times u$, full column rank unless stated otherwise), the
unknown parameters $\vc x$, the cofactor matrix $\mat\Sigma$ and the variance of unit weight
$\sigma_0^2$. The best linear unbiased estimate and its covariance are
\begin{equation}\label{eq:lsq}
\hat{\vc x}=\mat N^{-1}\vc n,\quad \mat N=\mat A\T\mat P\mat A,\quad \vc n=\mat A\T\mat P\vc l,\quad
\mat P=\mat\Sigma^{-1},\qquad
\hat\sigma_0^2=\frac{\hat{\vc v}\T\mat P\hat{\vc v}}{n-u},\quad
\mat\Sigma_{\hat x\hat x}=\hat\sigma_0^2\mat N^{-1},
\end{equation}
with the residuals $\hat{\vc v}=\vc l-\mat A\hat{\vc x}$. The residual is a projection,
$\hat{\vc v}=(\mat I-\mat H)\vc l$ with the hat matrix $\mat H=\mat A\mat N^{-1}\mat A\T\mat P$, and
the redundancy $n-u=\tr(\mat I-\mat H)$ is distributed over the observations by the diagonal of
$\mat R=\mat I-\mat H$, $\sum_i r_i=n-u$. Every variance estimator of Chapter~\ref{ch:vce} is a
specialisation of ``quadratic form over its expectation''.

\section{Whitening}\label{sec:whiten}

With a factor $\mat W$ of $\mat P=\mat W\T\mat W$ the model becomes
$\bar{\vc l}=\bar{\mat A}\vc x+\bar{\vc e}$ with $\bar{\mat A}=\mat W\mat A$, $\bar{\vc l}=\mat W\vc l$ and
$\mathrm C\{\bar{\vc e}\}=\sigma_0^2\mat I$, and \eqref{eq:lsq} is ordinary least squares. The
toolbox takes $\mat W=\mat L^{-1}$, the inverse of the lower Cholesky factor of
$\mat\Sigma=\mat L\mat L\T$. $\mat\Sigma$ is never formed for all observations: observations are
grouped into \emph{arcs}, independent of each other, and every arc is whitened by its own law
(Chapter~\ref{ch:stochastic}). All quantities below --- normal equations, residuals,
redundancies, test statistics --- are expressed in the whitened frame with a-priori
$\sigma_0=1$; ``residual'' means $\bar{\vc v}=\mat W\hat{\vc v}$ unless said otherwise.

\section{Two ways to accumulate and solve}\label{sec:backends}

\paragraph{Square-root information (QR, the default).} The upper-triangular factor $\mat R$ of
$[\bar{\mat A}\ \bar{\vc l}]$ is updated batch by batch by Householder reduction,
$\mat R\leftarrow\operatorname{triu}\operatorname{qr}\bigl([\mat R;\ \bar{\mat A}_k\ \bar{\vc l}_k]\bigr)$, and
truncated to $u+1$ rows. With
\[
\mat R=\begin{bmatrix}\mat R_1 & \vc r\\ \vc 0\T & \rho\end{bmatrix},\qquad
\hat{\vc x}=\mat R_1^{-1}\vc r,\qquad \bar{\vc v}\T\bar{\vc v}=\rho^2,\qquad
\mat\Sigma_{\hat x\hat x}=\hat\sigma_0^2(\mat R_1\T\mat R_1)^{-1},
\]
so the residual sum of squares is obtained without subtraction, and the condition number is
that of $\bar{\mat A}$, not its square. The rank decision is a pivot test: the system is
refused when $\min_k|R_{1,kk}|\le 10^3\,\varepsilon\,\max_k|R_{1,kk}|$.

\paragraph{Normal equations.} $\mat N=\sum_k\bar{\mat A}_k\T\bar{\mat A}_k$,
$\vc n=\sum_k\bar{\mat A}_k\T\bar{\vc l}_k$ and $\bar{\vc l}\T\bar{\vc l}$ are summed, and
$\mat N=\mat U\T\mat U$ is factorised by Cholesky. Partial systems add, which is the exchange
format of the field, but $\kappa(\mat N)=\kappa(\bar{\mat A})^2$ and
\begin{equation}\label{eq:cancel}
\bar{\vc v}\T\bar{\vc v}=\bar{\vc l}\T\bar{\vc l}-\vc n\T\hat{\vc x}
\end{equation}
cancels catastrophically when the fit is very good relative to $\|\bar{\vc l}\|$. Where the
cancellation exceeds the attainable precision the value is not trusted: the residual sum is
recomputed from the residuals in a second pass over the data.

\section{Block-arrow systems and the elimination of local parameters}\label{sec:blockarrow}

The parameters split into \emph{common} parameters $\vc x_c$ ($n_c$), seen by every observation,
and \emph{local} parameters $\vc y_s$ ($n_{l,s}$) of system $s$ --- initial states, instrument
calibrations, empirical terms --- valid for one arc, one day or the whole system. For the rows of
system $s$
\begin{equation}\label{eq:blockarrow}
\bar{\vc l}_s=\bar{\mat A}_s\vc x_c+\bar{\mat B}_s\vc y_s+\bar{\vc e}_s .
\end{equation}
The normal matrix has block-arrow form. The local block of each system --- all its arc-, day-
and system-scoped parameters together --- is eliminated by the Schur complement
\begin{equation}\label{eq:schur}
\mat N_{cc}^{\mathrm{red}}=\mat N_{cc}-\mat N_{lc}\T\mat N_{ll}^{-1}\mat N_{lc},\qquad
\vc n_c^{\mathrm{red}}=\vc n_c-\mat N_{lc}\T\mat N_{ll}^{-1}\vc n_l,\qquad
(\bar{\vc l}\T\bar{\vc l})^{\mathrm{red}}=\bar{\vc l}\T\bar{\vc l}-\vc n_l\T\mat N_{ll}^{-1}\vc n_l ,
\end{equation}
and the reduced systems of all systems add. On the square-root backend the same reduction is a
re-triangularisation with the local columns first,
\[
\begin{bmatrix}\mat R_{11}&\mat R_{12}&\vc r_1\\ & \mat R_{22}&\vc r_2\\ && \rho\end{bmatrix},
\]
whose trailing rows $[\mat R_{22}\ \vc r_2;\ \vc 0\T\ \rho]$ are the reduced contribution, and
$\mat R_{11},\mat R_{12},\vc r_1$ give the local solution and cofactor without forming
$\mat N_{ll}$. Each system contributes $m_s-n_{l,s}$ degrees of freedom. A local block that fails
the pivot test is rank-deficient: with $\mat R_{11}=\mat U\mat S\mat V\T$ and $k$ the singular values
above the rank tolerance, the locals live on $\mat V_1$ (the minimum-norm $\hat{\vc y}$, with
$\mat N_{ll}^{-1}$ read as the pseudo-inverse $\mat V_1\mat S_1^{-2}\mat V_1\T$ below), the rows
$\mat U_2\T[\mat R_{12}\ \vc r_1]$ join the reduced contribution, and the system contributes
$m_s-k$ degrees of freedom. Forming $\mat N_{ll}$ and cutting its eigenvalues instead loses a
weak but determined direction together with the null one.

Back-substitution recovers
\begin{equation}\label{eq:backsub}
\hat{\vc y}_s=\mat N_{ll}^{-1}(\vc n_l-\mat N_{lc}\hat{\vc x}_c),\qquad
\mat\Sigma_{\hat y\hat y}=\hat\sigma_0^2\mat N_{ll}^{-1}+\mat G\mat\Sigma_{\hat x\hat x}\mat G\T,\quad
\mat G=\mat N_{ll}^{-1}\mat N_{lc}.
\end{equation}
Every diagnostic that involves the hat matrix includes the local block. With
$\vc a=[\vc a_c;\vc b]$ a row in the full parameter space and $\vc z=\vc a_c-\mat G\T\vc b$,
\begin{equation}\label{eq:leverage}
\vc a\T\mat Q_{xx}\vc a=\bigl\|\mat R_c^{-\mathsf T}\vc z\bigr\|^2+\bigl\|\mat H_s\vc b\bigr\|^2,\qquad
\mat H_s\T\mat H_s=\mat N_{ll}^{-1},
\end{equation}
where $\mat R_c$ is the factor of the reduced common system; the full $u\times u$ cofactor matrix
is never formed.

\section{Prior information and regularisation}\label{sec:priors}

A prior is a group $k$ of $p_k$ pseudo-observations
\begin{equation}\label{eq:prior}
\mat L_k\vc x=\vc x_{0,k}+\vc e_k,\qquad \mathrm C\{\vc e_k\}=\sigma_k^2\mat I ,
\end{equation}
written in the columns of the parameters. In the whitened frame it contributes the rows
$\mat L_k/\sigma_k$ with observations $\vc x_{0,k}/\sigma_k$; this is Tikhonov regularisation with
$\lambda_k=\sigma_0^2/\sigma_k^2$, $\mat N_\lambda=\mat N+\sum_k\mat L_k\T\mat L_k/\sigma_k^2$.
Relative weights inside a group are carried by the rows of $\mat L_k$; $\sigma_k$ is one common
factor and can be estimated (\S\ref{sec:priorvce}).

Priors on common parameters are stacked into one operator $\mat K=[\mat L_k/\sigma_k]_k$ and
applied at the solve: on the square-root backend by appending its rows to the factor, on the
normal backend as $\mat N+\mat K\T\mat K$. A prior on local parameters is an observation stream of
its own with $\bar{\mat A}=\vc 0$, $\bar{\mat B}=\mat L/\sigma$ and $\bar{\vc l}=\vc x_0/\sigma$, so it
enters the elimination \eqref{eq:schur} like data. Either way the prior rows count as
observations: the redundancy is $n+\sum_k p_k-u$, and the share the data overrule is the
redundancy of the group,
\begin{equation}\label{eq:priorred}
r_k=p_k-\tr\bigl(\mat L_k\mat N_\lambda^{-1}\mat L_k\T\bigr)/\sigma_k^2 .
\end{equation}

\subsection{The datum of a free network}\label{sec:datum}
A rank defect --- a datum that no observation fixes --- is removed by a prior on the null space.
Let $\mat K_g$ be the rows of group $g$. The data are \emph{blind} to the group when
\[
\|\mat N\mat K_g\T\|_F\le\epsilon\,\bigl\|\,|\mat N|\,|\mat K_g\T|\,\bigr\|_F
\]
(elementwise absolute values: the product against its own rounding bound, which does not change
with the units of the parameters; on the square-root backend $\mat N$ is replaced by the factor).
A blind group is a \emph{datum} when it is also orthogonal to the rows $\mat K_o$ of every group
the data see, $\|\mat K_o\mat K_g\T\|_F\le\epsilon\,\|\mat K_o\|_F\|\mat K_g\|_F$; a datum may be
split over several groups. With $\mat K_d$ the rows
of all datum groups and $\mat K_i$ the rest, $\mat N_i=\mat N+\mat K_i\T\mat K_i$. Because the range
of the data and the null space are orthogonal, the estimate $\hat{\vc x}$ is the
minimum-constraint solution for any datum weight, and its covariance is that of the
minimum-constraint problem,
\begin{equation}\label{eq:mincon}
\mat\Sigma_{\hat x\hat x}=\hat\sigma_0^2\,\mat N_i^{+},
\end{equation}
evaluated stably as $\hat\sigma_0^2(\mat R_m^{-1}\mat R_m^{-\mathsf T}-\mat D_m\mat D_m\T)$ with
$\mat R_m\T\mat R_m=\mat N_i+c^2\mat K_d\T\mat K_d$, $\mat D_m=\mat K_d\T(\mat K_d\mat K_d\T)^{-1}/c$ and a
balancing constant $c$. A group only partly in the null space, or one that overlaps another group,
is information: its estimate moves with its $\sigma$, and it keeps the covariance
$\hat\sigma_0^2(\mat N+\mat K\T\mat K)^{-1}$.

\section{Nonlinear models: damped Gauss--Newton}\label{sec:gn}

For $\vc l=f(\vc x)+\vc e$ the objective, prior included, is
\begin{equation}\label{eq:gnobj}
F(\vc x)=\bigl\|\mat W(\vc l-f(\vc x))\bigr\|^2+\sum_k\bigl\|(\mat L_k\vc x-\vc x_{0,k})/\sigma_k\bigr\|^2 .
\end{equation}
At the iterate $\vc x$ the Jacobian $\mat J$ is accumulated once; that accumulation is both the
linearisation and the evaluation of $F(\vc x)$. The step $\Delta\vc x$ is one least-squares solve
with two kinds of pseudo-rows on the common block:
\begin{equation}\label{eq:gnstep}
\mat K\,\Delta\vc x\approx\vc x_0-\mat K\vc x,\qquad
\sqrt{\lambda d_j}\,\Delta x_j\approx0,\quad \vc d=\diag(\mat J\T\mat P\mat J).
\end{equation}
The first form makes the iteration converge to the minimiser of \eqref{eq:gnobj}; the second is
Marquardt damping of the parameters themselves. A step is accepted only if
$F(\vc x+\Delta\vc x)\le F(\vc x)$; a rejected step multiplies $\lambda$ by ten (from zero to
$10^{-3}$) and is re-solved on the same accumulation, an accepted one divides it by ten (below
$10^{-12}$ it is set to zero). The iteration stops when $\max_j|\Delta x_j|/\sigma_{x_j}$ or the
relative change of $F$ falls below its tolerance. The covariance is that of the undamped
linearised problem at the final point, so it does not depend on the damping history; a run that
exhausts its budget returns the best accepted point.

\section{What the estimates mean}\label{sec:meaning}

$\mat\Sigma_{\hat x\hat x}=\hat\sigma_0^2\mat N^{-1}$ is the \emph{formal} covariance: correct if the
functional model is complete, the stochastic model is right and the residuals are white.
Chapter~\ref{ch:vce} makes the second condition true as far as the data allow;
Chapter~\ref{ch:quality} tests the first and the third.
